% Direct-authority diplomatic French for printed p.198.
\clearpage
\SourcePage{198}

\LayerTarget{tohoku:construction:direct-image-g-sheaf:p1}
entendu, si $A$ est un $G$-faisceau de groupes ou d'anneaux etc., il en est de
même de $f_*(A)$.
\LayerTarget{tohoku:definition:invariant-direct-image}
Si $A$ est un $G$-faisceau, nous désignerons par $A^G$ ou $f_*^G(A)$ le
faisceau $\mathbf{\Gamma}^G(f_*(A))$ des invariants du $G$-faisceau $f_*(A)$ sur $Y$,
c'est-à-dire le faisceau dont les sections sur un ouvert $U\subset Y$ sont les
sections de $A$ sur $f^{-1}(U)$ \textit{invariantes sous $G$} ; évidemment, si $A$ est
un $G$-faisceau de groupes, ou anneaux etc, il en est encore de même de
$A^G=f_*^G(A)$.
\LayerTarget{tohoku:construction:direct-image-g-o-module}
D'ailleurs $f_*^G$ peut être considéré comme un \textit{foncteur} covariant
défini sur la catégorie des $G$-faisceaux sur $X$, à valeurs dans la catégorie
des faisceaux sur $Y$ (considérés comme $G$-faisceaux où $G$ opère
trivialement). Si ${\mathbf O}$ est un $G$-faisceau d'anneaux et si on désigne par ${\mathbf O}'$
le faisceaux $f_*^G({\mathbf O})$, alors pour tout $G$-${\mathbf O}$-Module $A$ sur $X$, $f_*(A)$
est un $G$-${\mathbf O}'$-Module sur $Y$. Donc $f_*$ est un \textit{foncteur} covariant
de ${\mathbf C}^{{\mathbf O}(G)}$ dans ${\mathbf C}^{{\mathbf O}'(G)}$, d'ailleurs additif et exact à gauche comme
tout le monde ; $f_*^G$ apparaît alors comme le foncteur composé
$\mathbf{\Gamma}^Gf_*$ de ${\mathbf C}^{{\mathbf O}(G)}$ dans la catégorie ${\mathbf C}^{{\mathbf O}'}$ des ${\mathbf O}'$-Modules sur
$Y$, $\mathbf{\Gamma}^G$ désignant le foncteur
${\mathbf C}^{{\mathbf O}'(G)}\to {\mathbf C}^{{\mathbf O}'}$, associant à un $G$-${\mathbf O}'$-faisceau $B$ sur
$Y$ le faisceau $B^G$ des germes de sections invariantes sous $G$. Supposons
maintenant pour simplifier que ${\mathbf O}$ donc aussi ${\mathbf O}'$ est le faisceau constant
$\mathbf Z$ des entiers (ce qui revient à ne pas considérer de faisceau
d'anneaux du tout), on a alors :

\LayerTarget{tohoku:proposition:5.1.3}
\textsc{Proposition 5.1.3} \textit{$f_*$ transforme objets injectifs
de ${\mathbf C}^{X(G)}$ en objets injectifs de ${\mathbf C}^{Y(G)}$.}

\LayerTarget{tohoku:proof:5.1.3}
En effet, il suffit de le voir pour un $G$-faisceau abélien $A$ défini par une
famille $(A_y)$ de $G_y$-modules injectifs (prop. 5.1.2). Mais reprenant les
notations de la construction qui a précédé la prop. 5.1.2, on voit que
$f_*(A)$ est le faisceau produit défini par les groupes $\overline A_y=\displaystyle\prod_{x\in f^{-1}(y)}A_x$,
et la structure de $G$-faisceau sur $f_*(A)$ est celle définie par la structure
de $G$-module de $\overline A_y$. Or ce dernier étant obtenu à partir du
$G_y$-module injectif $A_y$ par extension contravariante des scalaires, est
évidemment lui-même injectif [6, II, prop. 6.1.a]. Donc $f_*(A)$ est injectif
en vertu de prop. 3.1.2.

\LayerTarget{tohoku:corollary:5.1.3}
\textsc{Corollaire.} \textit{Si $A$ est un $G$-faisceau abélien
injectif, alors $\Gamma(A)$ est un $G$-module injectif, et $A^G$ est un
faisceau flasque sur $Y$.}

\LayerTarget{tohoku:proof:corollary-5.1.3}
En effet, on a $\Gamma(A)=\Gamma(f_*(A))$, $f_*^G(A)=\mathbf{\Gamma}^G(f_*(A))$, et il
suffit alors d'appliquer le lemme de 4.3 resp. la prop. 4.1.3 (avec
${\mathrm O}=\mathbf Z(G)$, ${\mathrm O}$ le faisceau d'anneaux constant défini par ${\mathrm O}$).

\LayerTarget{tohoku:construction:inverse-image-g-sheaf}
\textit{Images inverses et images directes.} Soit $B$ un faisceau
d'ensembles sur $Y$ (sans opérateurs), alors \textit{son image réciproque
$f^{-1}(A)$} (envisagée dans 3.2) \textit{peut être regardée comme un
$G$-faisceau} de façon naturelle, pour des raisons évidentes de ``transport
de structure''. Une section de $B$ sur un ouvert $U$ définit par image
réciproque une section de $f^{-1}(B)$ sur $f^{-1}(U)$ \textit{invariante sous
$G$}, i.e. section de $f_*^G(f^{-1}(B))$, d'où un homomorphisme naturel
$B\to f_*^G(f^{-1}(B))$, dont on vérifie aussitôt que c'est là un
isomorphisme :

\LayerTarget{tohoku:equation:5.1.1}
\[
f_*^G(f^{-1}(B))=B.
\tag{5.1.1}
\]
\LayerTarget{tohoku:construction:canonical-map-inverse-direct:p1}
Inversement, partons d'un $G$-faisceau $A$ sur $X$, considérons
$f^{-1}(f_*^G(B))$ : une section de ce faisceau sur un ouvert $V$ consiste en
la donnée d'une fonction

\TohokuContinuation{199}
